\chapter{Valeur absolue et fonction carré}
\section*{Exercice 14.1 -- Distances}
L'égalité donne $x=-3$ ou $x=5$. L'inégalité décrit tous les points à distance au plus $4$ du centre $1$, donc $[-3,5]$.

\section*{Exercice 14.2 -- Carrés}
La première donne $]-4,4[$. La seconde donne $]-\infty,-\sqrt5]\cup[\sqrt5,+\infty[$. Les bornes sont incluses seulement pour l'inégalité large.

\section*{Exercice 14.3 -- Comparer}
La différence $x^2-x=x(x-1)$ est négative pour $0<x<1$, donc $0{,}7^2<0{,}7$. Elle est positive pour $x>1$, donc $1{,}3^2>1{,}3$.

\section*{Exercice 14.4 -- Parabole déplacée}
Le carré est minimal lorsque $x+2=0$, donc en $x=-2$, avec valeur $-3$. La fonction décroît sur $]-\infty,-2]$ et croît sur $[-2,+\infty[$, par translation de la fonction carré.

\section*{Exercice 14.5 -- Solutions et domaine}
On obtient $(x+2)^2=9$, donc $x+2=-3$ ou $3$ : les solutions réelles sont $-5$ et $1$. Sur $[0,+\infty[$, seule $1$ est admissible.

\section*{Exercice 14.6 -- Un contre-exemple}
$-3<-2$, mais $9>4$. Avec $0\leq a<b$, on a $b^2-a^2=(b-a)(b+a)>0$. La restriction aux positifs rend l'affirmation correcte.
